\documentclass[10pt,letterpaper]{article}
\usepackage{vigo-developer-guide}

\begin{document}

\guideheader

\section{Quickstart}

This printable reference takes one GTFS feed and one OSM extract through Build and the first Route. The example uses Boston coordinates. Replace \code{YYYY-MM-DD} with an exact local service date covered by your feed.

\subsection{Install VIGO}

VIGO 0.4.2 requires Node.js 24.18 or newer, npm 11.6 or newer, and the pinned Rust toolchain for source builds. Native targets are macOS Apple Silicon/Intel, Linux ARM64/x64 with glibc, and Windows x64. Install the matching runtime; complete City data can move between these targets.

\begin{lstlisting}[language=bash]
git clone https://github.com/hytangs/vigo.git
cd vigo
npm ci
npm run build
npm link
\end{lstlisting}

\subsection{Build Boston}

\begin{lstlisting}[language=bash]
vigo build \
  --gtfs ./mbta.zip \
  --gtfs-scope mbta \
  --osm ./massachusetts.osm.pbf \
  --output ./boston
\end{lstlisting}

Build writes one complete City directory. The GTFS and OSM coverage must overlap.

\subsection{Run the first Route}

Save this as \code{route.json}:

\begin{lstlisting}[language=json]
{
  "origin": {"coordinate": [-71.11902, 42.37334]},
  "destination": {"coordinate": [-71.08337, 42.32978]}
}
\end{lstlisting}

\begin{lstlisting}[language=bash]
vigo route \
  --city ./boston \
  --request ./route.json \
  --time 11:04 \
  --service-date YYYY-MM-DD \
  --output ./route-result.json
\end{lstlisting}

An abbreviated Result has this shape (illustrative values):

\begin{lstlisting}[language=json]
{
  "kind": "route",
  "status": "ready",
  "result": {
    "durationMinutes": 33.517,
    "transfers": 2,
    "legs": ["walk", "ride", "walk", "ride", "..."]
  },
  "timing": {"computeMs": 3.301}
}
\end{lstlisting}

Values depend on the supplied City and request. Continue with Matrix and Reach below. The documentation index at \path{docs/README.md} links worked baseline/alternative workflows, Result interpretation, and Studio network investigations.

\section{Core model}

VIGO has four first-class concepts:

\begin{center}
\begin{tabularx}{\linewidth}{P{1.0in} Y}
\toprule
\textbf{Concept} & \textbf{Meaning} \\
\midrule
City & An immutable compiled mobility model made from GTFS and OSM. \\
Scenario & Optional immutable changes tied to one City revision. Baseline work uses the City directly. \\
Query & One explicit Route, Matrix, or Reach request. \\
Result & The answer together with status, normalized input, warnings, timing, and City identity. \\
\bottomrule
\end{tabularx}
\end{center}

\begin{center}
\sffamily\bfseries\color{VigoTeal}
City $\longrightarrow$ Scenario? $\longrightarrow$ Route $\mid$ Matrix $\mid$ Reach $\longrightarrow$ Result
\end{center}

Every computation returns a Result, not a naked travel time.

\subsection{Three Queries}

\begin{center}
\begin{tabularx}{\linewidth}{P{0.9in} Y Y}
\toprule
\textbf{Query} & \textbf{Question} & \textbf{Principal output} \\
\midrule
Route & How can one journey move through the City? & Ordered walk, ride, or drive legs. \\
Matrix & How long does travel take between sets of places? & One scalar row per requested pair. \\
Reach & Where can the modeled network reach within stated time limits? & A travel-time surface and contours. \\
\bottomrule
\end{tabularx}
\end{center}

Compare is not a fourth Query. It operates on two compatible saved Results and does not rerun either computation:

\begin{center}
\sffamily\bfseries\color{VigoNavy}Compare($Result_a$, $Result_b$) $\longrightarrow$ Comparison Result
\end{center}

\subsection{Interfaces}

\begin{center}
\begin{tabularx}{\linewidth}{P{1.0in} P{2.0in} Y}
\toprule
\textbf{Interface} & \textbf{Role} & \textbf{Use it for} \\
\midrule
CLI & Shell and automation & Builds, scripts, repeatable runs, and CI. \\
Studio & Visual & City exploration, route inspection, playback, and interactive analysis. \\
\bottomrule
\end{tabularx}
\end{center}

The interfaces share nouns and result meaning. Their presentation can differ without creating another product model.

\section{Cities}

\subsection{Inputs and Build}

A City needs one or more static GTFS ZIP files, one overlapping OSM \code{.osm.pbf}, a new output directory, and enough storage for compilation. Repeat \field{--gtfs} for multiple feeds. Give each feed a stable scope when identifiers may overlap.

\begin{lstlisting}[language=bash]
vigo build \
  --gtfs ./operator-a.zip --gtfs-scope operator-a \
  --gtfs ./operator-b.zip --gtfs-scope operator-b \
  --osm ./metro-region.osm.pbf \
  --output ./metro-region
\end{lstlisting}

VIGO publishes the output only after required timetable and street files are complete. An existing output directory is left alone unless \field{--replace} is supplied.

Public pedestrian access is the default. Use \field{--private-access=endpoints} when travelers may use internal roads at their own endpoints. This includes mapped private access within the same walking budget. Private streets remain unavailable for through travel and transit transfers. This City-wide assumption does not infer ownership or gate hours. Rebuild older Cities for the current street-store v4 permission rules.

\subsection{Inspect and identify}

\begin{lstlisting}[language=bash]
vigo inspect --city ./metro-region
\end{lstlisting}

Inspect reports the City name, revision identity, build time, sources, and network counts. Its compiled network is immutable; replacement inputs produce a new revision. Reopen the complete City directory without re-importing raw sources. Street indexes, station-access state, and retained timetable snapshots are reusable across processes. A new active-service pattern may require one timetable preparation. Move or archive the whole directory to retain its prepared files.

Points are either GTFS stop IDs or coordinates. Coordinates always use \textbf{[longitude, latitude]}. Coordinate routing connects through the OSM street graph.

\section{Results and failures}

\subsection{Common Result fields}

\begin{center}
\begin{tabularx}{\linewidth}{P{1.25in} Y}
\toprule
\textbf{Field} & \textbf{Meaning} \\
\midrule
\field{kind} & \code{route}, \code{matrix}, \code{reach}, or \code{comparison}. \\
\field{status} & \code{ready} or \code{blocked}. \\
\field{query} & Normalized request summary; retain the original JSON and command for complete inputs. \\
Query output & Route uses \field{result}; Matrix uses \field{rows}; Reach uses \field{surface} and \field{contours}; Compare uses \field{change}. \\
\field{warnings} & Reported qualifications; also inspect diagnostics and individual legs. \\
\field{timing} & Separate available durations such as open, compute, contour, and end to end. \\
\field{city} & Query City identity; Compare instead records both inputs under \field{cities}. \\
\bottomrule
\end{tabularx}
\end{center}

Fields vary by Result family. Compare returns its own envelope, not copies of both inputs. An empty \field{warnings} array does not establish complete source coverage. In Transit Route, inspect \field{result.diagnostics} for admitted realtime updates, scheduled fallback, and access qualifications; a requested realtime mode alone does not prove that predictions were applied.

\subsection{Status and exceptions}

\begin{center}
\begin{tabularx}{\linewidth}{P{1.35in} Y}
\toprule
\textbf{Name} & \textbf{Meaning} \\
\midrule
\code{ready} & Computation completed; Matrix may still contain blocked rows, and Reach may contain unreachable cells. \\
\code{blocked} & No usable Route journey was returned under the supplied inputs and constraints. \\
Invalid input & The command exits nonzero with an explanation on standard error. \\
Unsupported request & The selected City or Scenario cannot execute the request; the command exits nonzero. \\
Execution failure & Build or Query execution failed; the command exits nonzero. \\
\bottomrule
\end{tabularx}
\end{center}

A blocked Route Result exits zero. Invalid input and execution failures exit 2. Count Matrix and CSV Route rows by their own status; missing durations are not zero-minute journeys.

JSON query commands print the complete Result to stdout; \field{--output PATH} also saves it to a file. A CSV Route batch writes rows to \field{--output} and prints its JSON summary, including detailed plans, to stdout. Redirect stdout separately to retain that summary.

\subsection{Keep a reproducible run}

Retain the complete City, \code{vigo inspect} output, runtime version and capabilities, original request, exact command, full Result, and stderr together. Keep supplied realtime or traffic snapshots and any comparison inputs. Preserve source archives and acquisition notes when a rebuild must be reproducible. A normalized \field{query} or provenance key does not archive all these inputs.

Live admission checks a captured current clock. A saved snapshot can expire before a later query; the public CLI provides no historical replay clock. Preserve the original Result for incident review. See \path{docs/reference/results.md} for the retention checklist.

\section{Route}

\subsection{What it answers}

Route finds and explains one journey through ordered points. Depart-at, arrive-by, transit, walk, drive, departure windows, waypoints, and batch requests remain Route options.

\subsection{Minimal examples}

\begin{lstlisting}[language=bash]
vigo route --city ./boston --request ./route.json \
  --time 11:04 --service-date YYYY-MM-DD
\end{lstlisting}

\subsection{Inputs and options}

\begin{longtable}{P{1.72in} P{1.04in} P{3.54in}}
\toprule
\textbf{Input or option} & \textbf{Default} & \textbf{Meaning} \\
\midrule
\endhead
\field{origin} & required & Exact stop ID or point object with \field{coordinate}. \\
\field{destination} & required & Exact stop ID or point object with \field{coordinate}. \\
\field{waypoints} & empty & Ordered intermediate points. \\
\field{--mode} & transit & \code{transit}, \code{walk}, or \code{drive}. \\
\field{--time} & 08:00 & Local departure or arrival time. \\
\field{--time-preference} & depart & \code{depart} or \code{arrive}. \\
\field{--service-date} & required & Exact GTFS date in \code{YYYY-MM-DD} form. \\
\field{--max-walk} & 1.2 & Transit access and egress distance in km. \\
\field{--max-transfers} & unlimited & At most 0--31 vehicle changes; unavailable with transit waypoints. \\
\field{--horizon} & 480 & Timetable search horizon, 1--2880 minutes. \\
\field{--departure-window} & 0 & Depart-at profile of plus or minus $N$ integral minutes, $N$ from 0 to 30. \\
\field{--objective} & earliest arrival & Earliest arrival, then fewer boardings, less walking, and stable order. \\
\bottomrule
\end{longtable}

\subsection{Result}

Route returns departure and arrival times, total duration, boarding count, ordered legs, geometry, warnings, and timing. A leg identifies its type and endpoints. Ride legs add route and trip detail; street legs add distance and the selected OSM path.

\subsection{Important behavior}

Transit requires an exact service date. VIGO applies \code{calendar.txt} and \code{calendar\_dates.txt}. It does not substitute a generic weekday outside the feed's active service.

Arrive-by uses directed OSM access and egress, an exact reverse timetable scan for the latest feasible physical departure, and forward reconstruction of the selected journey in the same computation.

\begin{lstlisting}[language=bash]
vigo route --city ./boston --request ./route.json \
  --time 09:00 --time-preference arrive \
  --service-date YYYY-MM-DD
\end{lstlisting}

Transit Route requires a vehicle boarding by default. Set \field{requireTransitRide: false} in JSON to compare transit with walking, or use Walk mode for a walking journey. A valid request with no journey returns \code{blocked}.

\subsection{Advanced Route}

Departure windows sample each integral minute from the requested departure minus $N$ through the requested departure plus $N$. They are depart-at only.

Walking and driving require coordinates. Drive uses free-flow street weights unless fresh supplied traffic is applied in realtime mode. Waypoints are visited in listed order; VIGO does not reorder them.

For batch work, pass CSV through \field{--input} and write the returned rows through \field{--output}:

\begin{lstlisting}[language=bash]
vigo route --city ./boston \
  --input ./routes.csv --output ./route-results.csv \
  --time 08:00 --service-date YYYY-MM-DD
\end{lstlisting}

\section{Matrix}

\subsection{What it answers}

Matrix computes scalar travel time between sets of origins and destinations. The same Query covers one-to-many, many-to-one, and many-to-many work.

\noindent\begin{minipage}{\linewidth}
\subsection{Minimal examples}

Save \code{matrix.json}:

\begin{lstlisting}[language=json]
{
  "origins": [
    {"id":"harvard", "point":{"coordinate":[-71.11902,42.37334]}},
    {"id":"central", "point":{"coordinate":[-71.10350,42.36540]}}
  ],
  "destinations": [
    {"id":"back_bay", "point":{"coordinate":[-71.07540,42.34730]}},
    {"id":"nubian", "point":{"coordinate":[-71.08337,42.32978]}}
  ]
}
\end{lstlisting}

\begin{lstlisting}[language=bash]
vigo matrix --city ./boston --request ./matrix.json \
  --time 08:00 --service-date YYYY-MM-DD \
  --horizon 180 --output ./matrix-result.json
\end{lstlisting}

\end{minipage}

\subsection{Inputs and options}

\begin{longtable}{P{1.64in} P{1.10in} P{3.56in}}
\toprule
\textbf{Input or option} & \textbf{Default} & \textbf{Meaning} \\
\midrule
\endhead
\field{origins} & required & Array of points or \code{\{id, point\}} objects. \\
\field{destinations} & required & Array of points or \code{\{id, point\}} objects. \\
\field{--mode} & transit & \code{transit}, \code{walk}, or \code{drive}. \\
\field{--time} & 08:00 & Local departure time or arrival deadline; select with \field{--time-preference}. \\
\field{--service-date} & required & Exact timetable date. \\
\field{--max-walk} & 1.2 & Transit access and egress distance in km. \\
\field{--horizon} & 480 & Timetable search horizon from 1 to 2880 minutes. \\
\field{walkSpeedKph} & 4.8 & JSON field used for Walk Matrix. \\
\field{maxDistanceKm} & none & Optional street-matrix distance limit. \\
\bottomrule
\end{longtable}

\subsection{Result}

Matrix preserves caller IDs and returns one row per requested pair in origin-major order.

Coordinate transit Matrix projects endpoints and optionally evaluates direct walking in Rust. Public street attachments may be shared within the same request; private endpoint access retains its own distance boundary. \field{disableCache: true} disables cross-request access/path reuse while keeping the prepared City.

Matrix returns scalar rows by default. Set \field{includeJourneys: true} to include actual arrival, walking, waiting, transfers, and timed legs. Add \field{includeGeometry: true} for the selected paths, without repeating timetable searches. Count \code{ready} and \code{blocked} rows separately.

\subsection{Limits and important behavior}

One request accepts at most 100,000 pairs and 16 MiB of JSON. Transit Matrix supports depart-at and arrive-by. Scalar searches share a forward scan per unique origin for depart-at, or a reverse scan per unique destination for arrive-by; requested journeys add shared journey rounds. Arrive-by scalar duration is the deadline minus latest departure, including destination waiting; a nested journey reports actual modeled arrival. Walk and Drive Matrix use the street graph. Supplied traffic applies only to Drive Matrix.

\section{Reach}

\subsection{What it answers}

Reach computes where the represented network can travel within stated time limits. It measures network reach. It does not count jobs, people, schools, healthcare, or other opportunities; those additions belong to a separate accessibility method.

\noindent\begin{minipage}{\linewidth}
\subsection{Minimal examples}

Save \code{reach.json}:

\begin{lstlisting}[language=json]
{
  "origin": {"coordinate": [-71.11902, 42.37334]},
  "cutoffsMinutes": [15, 30, 45],
  "extentRadiusKm": 6,
  "rasterSize": 48
}
\end{lstlisting}

\begin{lstlisting}[language=bash]
vigo reach --city ./boston --request ./reach.json \
  --time 08:00 --service-date YYYY-MM-DD \
  --output ./reach-result.json
\end{lstlisting}

\end{minipage}

\subsection{Inputs and options}

\begin{longtable}{P{1.62in} P{1.08in} P{3.60in}}
\toprule
\textbf{Input or option} & \textbf{Default} & \textbf{Meaning} \\
\midrule
\endhead
\field{origin} & required & Exact stop ID or point object with \field{coordinate}. \\
\field{--time} & 08:00 & Fixed departure time. \\
\field{--service-date} & required & Exact timetable date. \\
\field{--cutoffs} & 15,30,45,60 & Increasing limits from 5 to 240 minutes. \\
\field{--max-walk} & 1.2 & Transit access and terminal walking distance in km. \\
\field{--walk-speed} & 4.8 & Walking speed from 1 to 8 kph. \\
\field{--extent-radius} & 8 & Requested raster radius from 1 to 40 km. \\
\field{--raster-size} & 96 & Grid size: 48, 64, 96, 128, 192, 256, 384, 512, or 1024. \\
\bottomrule
\end{longtable}

\subsection{Result}

Reach returns a bounded travel-time surface and drawable contours. The surface retains its coordinate bounds. The extent radius defines the requested CLI raster, not a maximum trip distance. Cutoffs and walking budgets govern network reachability. Studio expands its displayed surface to the complete reached-network envelope. The raster size controls grid resolution. Retain both with the Result.

A 48 by 48 grid contains 2,304 cells. Surface coverage and contour availability are separate facts.

\section{Scenario and Compare}

\subsection{Scenario in 0.4.2}

A Scenario is an immutable set of supported changes tied to exactly one City revision. It does not edit the City. Time, mode, walking limits, cutoffs, and objective belong to the Query, not the Scenario.

The abstraction is stable; support in 0.4.2 is deliberately narrower:

\begin{center}
\begin{tabularx}{\linewidth}{P{1.42in} Z Z Z Y}
\toprule
\textbf{State} & \textbf{Route} & \textbf{Matrix} & \textbf{Reach} & \textbf{Available in 0.4.2} \\
\midrule
Baseline City & \yes & \yes & \yes & CLI; Studio offers Route and Reach \\
Planned service & \no & \no & \yes & CLI and Studio \\
Traffic & Drive & Drive & \no & CLI \\
Live transit & \yes & \no & \no & Studio; CLI with supplied snapshot \\
\bottomrule
\end{tabularx}
\end{center}

Planned service can add scheduled service or exclude selected routes for Reach. Traffic carries time-bounded directed road observations for Drive Route and Drive Matrix. Live transit Route processes supported Trip Updates against matched scheduled trips in Studio and the CLI. CLI callers supply \field{realtimeSnapshot} and select \code{--data-mode realtime}. The complete supplied snapshot is processed without trip-count caps; unreported trips retain scheduled times. Cancellations and skipped calls are supported; added trips without a scheduled instance remain unsupported. Vehicle Positions and Alerts do not change routing. Inspect realtime diagnostics for freshness exclusions and scheduled fallback. Transit Matrix and Reach remain scheduled.

Use \code{vigo capabilities} to inspect support. Supply planned changes in Reach's \field{scenario} object. Drive traffic uses a top-level \field{traffic} object and explicit realtime selection: Route accepts \code{--data-mode realtime}; Matrix requires \code{"routingDataMode": "realtime"} in JSON and has no \field{--data-mode} flag. Scheduled mode ignores supplied observations. Unsupported planned-service combinations are rejected.

\subsection{Compare saved Results}

Compare accepts two Results from the same Query family. Route comparison reads one journey from each input, not an entire CSV batch. Matrix joins origin/destination IDs, falling back to indexes when IDs are absent; unmatched rows are omitted. Reach requires identical bounds, dimensions, and value-array length.

\begin{lstlisting}[language=bash]
vigo compare \
  --before ./baseline-result.json \
  --after ./proposal-result.json \
  --output ./comparison.json
\end{lstlisting}

Changes are after minus before. Negative values mean faster modeled travel. Matrix and Reach means include only pairs or cells with finite values in both inputs; newly and no-longer reachable counts remain separate. A mean with no comparable values is \code{null}.

Compare does not verify every City, date, endpoint, or option identity. Align these explicitly, and keep Matrix IDs unique and stable. Use one City revision to isolate a planned Scenario; comparing source revisions studies a different change. The worked example is in \path{docs/guides/workflows.md}.

\section{Command-line reference}

\subsection{Commands}

\begin{longtable}{P{1.05in} P{2.10in} P{3.35in}}
\toprule
\textbf{Command} & \textbf{Required input} & \textbf{Output} \\
\midrule
\endhead
\code{build} & GTFS, OSM, output path & One complete City and build summary. \\
\code{capabilities} & none & Runtime versions and supported combinations. \\
\code{inspect} & City & City identity, sources, and counts. \\
\code{route} & City, request or CSV, date & Route Result JSON or batch CSV. \\
\code{matrix} & City, request, date & Matrix Result JSON. \\
\code{reach} & City, request, date & Reach Result JSON. \\
\code{compare} & before and after Results & Comparison Result without recomputation. \\
\bottomrule
\end{longtable}

\subsection{Shared query options}

\begin{longtable}{P{1.82in} P{4.68in}}
\toprule
\textbf{Option} & \textbf{Meaning} \\
\midrule
\endhead
\field{--city PATH} & Complete City directory. \\
\field{--request PATH} & One JSON request for Route, Matrix, or Reach; \code{-} reads stdin. \\
\field{--output PATH} & Save the complete Result. \\
\field{--mode VALUE} & \code{transit}, \code{walk}, or \code{drive}. \\
\field{--time HH:MM} & Local time, default 08:00. \\
\field{--time-preference} & \code{depart} or \code{arrive} for Route and Matrix. \\
\field{--service-date} & Required exact \code{YYYY-MM-DD} date, including CLI Walk and Drive. \\
\field{--max-walk KM} & Transit access and egress distance, default 1.2. \\
\field{--objective} & \code{earliest\_arrival} in 0.4.2. \\
\field{--departure-window MIN} & Route depart-at profile from 0 to 30 minutes. \\
\field{--horizon MIN} & Route and Matrix timetable horizon, default 480. \\
\field{--cutoffs MINUTES} & Reach limits separated by commas. \\
\field{--extent-radius KM} & Reach requested raster radius, default 8. \\
\field{--raster-size N} & Reach grid size from the supported set. \\
\field{--walk-speed KPH} & Reach walking speed, default 4.8. \\
\bottomrule
\end{longtable}

Successful commands exit zero and write JSON to standard output. Invalid input, an incomplete City, or execution failure exits nonzero and writes an explanation to standard error. A blocked Query remains a successful computation with a blocked Result.

JSON requests are limited to 16 MiB. Date and clock come from \field{--service-date} and \field{--time}; keep these flags with the request file. If a command fails, an existing output file may still contain an older Result.

\newpage
\section{Capability and compatibility}

\subsection{Interface support}

\begin{center}
\begin{tabularx}{\linewidth}{P{1.55in} Z Z Y}
\toprule
\textbf{Capability} & \textbf{Studio} & \textbf{CLI} & \textbf{Note} \\
\midrule
Build & Yes & Yes & Studio uses its project library; CLI builds movable Cities. \\
Route & Yes & Yes & Transit, walk, drive; depart-at and arrive-by. \\
Matrix & No & Yes & Scalar rows with optional transit journeys. \\
Reach & Yes & Yes & Surface and contours. \\
Planned service & Yes & Yes & Reach only. \\
Traffic state & No & Yes & Drive Route and Matrix. \\
Live transit & Yes & Yes & Route with a supplied snapshot. \\
Compare & Yes & Yes & Studio computes and compares Reach surfaces; CLI compares saved Results. \\
\bottomrule
\end{tabularx}
\end{center}

Run \code{vigo capabilities} to inspect support.

\subsection{Four versions}

\begin{lstlisting}[language=json]
{
  "productVersion": "0.4.2",
  "apiVersion": "1.0",
  "cityFormatVersion": 1,
  "resultSchemaVersion": 1
}
\end{lstlisting}

\field{productVersion} identifies the shipped release. \field{apiVersion} controls interface compatibility. \field{cityFormatVersion} identifies readable City data. \field{resultSchemaVersion} identifies saved Result structure.

\section{Troubleshooting}

\subsection{Coordinate query fails}

\textbf{Check.} Confirm longitude comes first and the coordinate lies inside the OSM extract. The City must contain \path{osm/street-index.sqlite}. Inspect street permissions and the selected distance limit.

\textbf{Fix.} Correct the coordinate or rebuild the City with an OSM extract that covers the endpoints. A working stop-to-stop Route does not prove coordinate connectivity.

\subsection{Date is rejected or returns blocked}

\textbf{Check.} Use \code{YYYY-MM-DD}; inspect the feed's service range and its \code{calendar.txt} and \code{calendar\_dates.txt} rules.

\textbf{Fix.} Select an active date or rebuild from the intended feed. VIGO does not shift the request to another date.

\subsection{Matrix has fewer usable pairs than expected}

\textbf{Check.} Count rows by status. Confirm the request limits, exact date, timetable horizon, street coverage, and walking distance.

\textbf{Fix.} Correct the limiting input. Run selected pairs as Route when you need their full journeys.

\subsection{Reach was interpreted as accessibility}

\textbf{Check.} A Reach Result contains transport impedance, not opportunity counts.

\textbf{Fix.} Join Matrix or Reach output to a separately documented opportunity dataset and method before reporting an accessibility measure.

\clearpage
\appendix

\section{Performance measurement}

\begin{center}
\begin{tabularx}{\linewidth}{P{1.05in} Y Y}
\toprule
\textbf{Term} & \textbf{Starts} & \textbf{Ends} \\
\midrule
Build & Invoke Build with local GTFS and OSM, with no existing output City & The command returns after publishing a complete City. \\
Raw to first answer & The same Build invocation & The first Query returns its complete Result. \\
Reopen & Start a new process for an existing City & Its first Query returns its complete Result. \\
Compute & Query execution begins & The Route, Matrix, or Reach answer exists. \\
End to end & Caller submits work & The complete Result returns. \\
\bottomrule
\end{tabularx}
\end{center}

A first Query may include Open work. Repeated work may use an already open City. A Result read from disk is not a new computation. Compare performance only with the same inputs, machine state, output detail, status accounting, and timing boundary.

Build prepares the street indexes, transfers, and station-access state. New Cities persist the Drive hierarchy and both metrics during Build; opening them loads those files. Older ephemeral Drive Cities retain their previous startup behavior until rebuilt from source. Required indexes must travel with the City. The first Query can additionally prepare its active-service timetable and selected ride geometry. Materialization-critical shape alignment and identifier arithmetic run in Rust; JavaScript still loads selected source rows and assembles itinerary objects. Downloads and installation are separate. In 0.4.2, \field{network.json} timing \field{totalMs} measures work inside the compiler; it excludes startup, final publication, and the first Query. Use a caller timer for complete Build and raw-to-first-answer measurements. Compiler stages can overlap and their times must not be summed.

For arrive-by Route, the reverse timetable scan is one part of Compute. Coordinate access, forward journey reconstruction, selected geometry, and Result creation also belong to full Route computation. Report the scan alone only when diagnosing that stage.

Large Matrix requests and fine Reach rasters are explicit caller choices. Record their dimensions. Selected ride shapes are aligned when first needed; opening a City does not prepare geometry for every active trip.

\section{City directory}

The main files are shown below. Keep all generated sidecar and native snapshot files with them:

\begin{lstlisting}
boston/
  network.json
  routing/
    project.sqlite
  osm/
    street-index.sqlite
\end{lstlisting}

Do not edit these files by hand. Move the directory as one unit. Rebuild when source GTFS or OSM changes.

\end{document}
